\chapter{Special Relativity and Spacetime Formulation}

\begin{remarkbox}
Classical field theory becomes structurally clear when it is written on
spacetime rather than on space plus time separately. This chapter introduces the
relativistic language needed for the rest of the course: Minkowski spacetime,
Lorentz transformations, four-vectors, tensors, proper time, causality,
covariant equations, and invariant actions.
\end{remarkbox}

\section{Minkowski Spacetime}

Special relativity is built on the idea that space and time are not independent
backgrounds. They form a single four-dimensional spacetime. A point of
spacetime is called an event and is labelled by coordinates
\[
    x^\mu = (x^0,x^1,x^2,x^3).
\]
In units where \(c=1\), one usually takes \(x^0=t\). If the speed of light is
kept explicit, one often writes \(x^0=ct\).

The geometric structure of flat relativistic spacetime is encoded in the
Minkowski metric. In these notes we use the mostly-minus convention
\[
    \eta_{\mu\nu}=\operatorname{diag}(+1,-1,-1,-1).
\]
The spacetime interval between neighbouring events is
\[
    ds^2 = \eta_{\mu\nu}dx^\mu dx^\nu
    = dt^2 - d\vect{x}^{\,2}
\]
in units \(c=1\). With \(c\) explicit, this becomes
\[
    ds^2 = c^2dt^2-d\vect{x}^{\,2}.
\]

The interval is not a Euclidean distance. Its sign has physical meaning. A
separation is timelike if \(ds^2>0\), lightlike if \(ds^2=0\), and spacelike if
\(ds^2<0\). Timelike separated events can lie on the worldline of a massive
particle. Lightlike separated events can be connected by a light ray. Spacelike
separated events cannot influence each other by any signal travelling at or
below the speed of light.

\begin{definitionbox}
\emph{Minkowski spacetime} is the flat spacetime of special relativity. It is
equipped with the metric \(\eta_{\mu\nu}\), which defines the invariant interval
\(ds^2=\eta_{\mu\nu}dx^\mu dx^\nu\).
\end{definitionbox}

The key difference between Euclidean geometry and Minkowski geometry is that
the metric has mixed signs. This is not a minor convention. It is the algebraic
origin of light cones, causal structure, relativistic wave equations, and the
classification of particles by mass.

\section{Lorentz Transformations}

A Lorentz transformation is a linear change of inertial coordinates that
preserves the Minkowski interval. If
\[
    x'^\mu = \Lambda^\mu{}_{\nu}x^\nu,
\]
then \(\Lambda^\mu{}_{\nu}\) is Lorentz if
\[
    \eta_{\rho\sigma}\Lambda^\rho{}_{\mu}\Lambda^\sigma{}_{\nu}
    = \eta_{\mu\nu}.
\]
Equivalently, in matrix notation,
\[
    \Lambda^T\eta\Lambda=\eta.
\]
This is the relativistic analogue of the condition \(R^TR=I\) for rotations in
Euclidean space, but with the Euclidean metric replaced by the Minkowski metric.

A standard boost in the \(x^1\)-direction has the form
\[
\begin{aligned}
    t' &= \gamma(t-vx^1),\\
    x'^1 &= \gamma(x^1-vt),\\
    x'^2 &= x^2,\\
    x'^3 &= x^3,
\end{aligned}
\]
where
\[
    \gamma = \frac{1}{\sqrt{1-v^2}}
\]
in units \(c=1\). This transformation mixes time and space. That mixing is why
relativistic field theory should be formulated on spacetime from the beginning.

\begin{examplebox}
For a boost in the \(x^1\)-direction, one can check directly that
\[
    t'^2-(x'^1)^2-(x'^2)^2-(x'^3)^2
    = t^2-(x^1)^2-(x^2)^2-(x^3)^2.
\]
The coordinates change, but the interval does not.
\end{examplebox}

The set of Lorentz transformations forms the Lorentz group, usually denoted
\(O(1,3)\). The transformations continuously connected to the identity and
preserving the direction of time form the proper orthochronous Lorentz group
\(SO^+(1,3)\). Most local relativistic field theories are built to transform
naturally under this group.

\section{Four-Vectors and Tensors}

A four-vector is an object whose components transform under Lorentz
transformations in the same way as spacetime coordinates:
\[
    V'^\mu = \Lambda^\mu{}_{\nu}V^\nu.
\]
Examples include the displacement \(dx^\mu\), the four-momentum \(p^\mu\), and
the four-current \(J^\mu\). A covector transforms with the inverse matrix, and
is usually written with a lower index:
\[
    W'_\mu = (\Lambda^{-1})^\nu{}_{\mu}W_\nu.
\]
The metric converts vectors into covectors:
\[
    V_\mu = \eta_{\mu\nu}V^\nu.
\]

The contraction of a vector and a covector is Lorentz invariant:
\[
    W_\mu V^\mu = W'_\mu V'^\mu.
\]
Similarly, the squared norm of a four-vector is
\[
    V_\mu V^\mu.
\]
For four-momentum,
\[
    p^\mu=(E,\vect{p})
\]
in units \(c=1\), the invariant relation is
\[
    p_\mu p^\mu = E^2-\vect{p}^{\,2}=m^2.
\]
This is the relativistic energy-momentum relation.

Tensors generalize vectors by carrying more indices. A rank-two tensor
\(T^{\mu\nu}\) transforms as
\[
    T'^{\mu\nu}
    = \Lambda^\mu{}_{\rho}\Lambda^\nu{}_{\sigma}T^{\rho\sigma}.
\]
The electromagnetic field strength \(F_{\mu\nu}\) is an antisymmetric rank-two
tensor. The stress-energy tensor \(T^{\mu\nu}\) is a rank-two tensor encoding
energy density, momentum density, energy flux, and stresses.

\begin{definitionbox}
A relativistic equation is \emph{covariant} if both sides transform in the same
way under Lorentz transformations. Covariance ensures that the equation has the
same geometric content in all inertial frames.
\end{definitionbox}

A useful rule is that a valid tensor equation must have the same free indices on
both sides. For example,
\[
    \partial_\mu F^{\mu\nu}=J^\nu
\]
is a valid vector equation, because the free index \(\nu\) appears on both
sides. The equation represents four component equations, but it is written in a
form that makes Lorentz covariance manifest.

\section{Proper Time, Worldlines, and Causal Structure}

The history of a point particle is a curve in spacetime, called a worldline:
\[
    x^\mu=x^\mu(\lambda),
\]
where \(\lambda\) is a parameter along the curve. For a massive particle, one
can use proper time \(\tau\) as the parameter. Proper time is defined by
\[
    d\tau^2 = ds^2
\]
for timelike curves in units \(c=1\). Thus
\[
    d\tau^2 = dt^2-d\vect{x}^{\,2}.
\]
If \(\vect{v}=d\vect{x}/dt\), then
\[
    d\tau = dt\sqrt{1-\vect{v}^{\,2}}
    = \frac{dt}{\gamma}.
\]
Proper time is the time measured by a clock moving along the worldline.

The four-velocity is defined as
\[
    u^\mu = \frac{dx^\mu}{d\tau}.
\]
It satisfies
\[
    u_\mu u^\mu = 1
\]
with our signature convention. The four-momentum is
\[
    p^\mu = m u^\mu,
\]
and therefore
\[
    p_\mu p^\mu=m^2.
\]

Causal structure is encoded by light cones. At each event, the light cone
separates possible future-directed timelike and lightlike directions from
spacelike directions. Relativistic field equations must respect this structure:
physical disturbances should propagate inside or on the light cone, not outside
it.

\begin{remarkbox}
The light cone is not an optical illustration added to spacetime. It is the
causal skeleton of relativistic physics. Hyperbolic field equations, retarded
Green's functions, and relativistic signal propagation are all organized around
this structure.
\end{remarkbox}

For field theory, worldlines are less fundamental than fields over spacetime,
but they remain useful. Sources, detectors, observers, and particles moving in
external fields are often represented by worldlines interacting with spacetime
fields.

\section{Relativistic Notation for Fields}

A relativistic field is written as a function of spacetime coordinates:
\[
    \phi(x) = \phi(x^0,x^1,x^2,x^3).
\]
For a scalar field, Lorentz transformations act by changing the argument:
\[
    \phi'(x')=\phi(x).
\]
Equivalently,
\[
    \phi'(x)=\phi(\Lambda^{-1}x).
\]
This means that the scalar value attached to a spacetime event is independent
of the coordinate system used to describe that event.

A vector field \(A_\mu(x)\) transforms both in its argument and in its index:
\[
    A'_\mu(x')=(\Lambda^{-1})^\nu{}_{\mu}A_\nu(x).
\]
A tensor field transforms with one Lorentz matrix for each index. Spinor fields
transform with spinor representations rather than ordinary vector
representations; this will be discussed later in the chapter on the Dirac
field.

The derivative operator is written as
\[
    \partial_\mu = \frac{\partial}{\partial x^\mu}.
\]
Raising the index gives
\[
    \partial^\mu=\eta^{\mu\nu}\partial_\nu.
\]
The d'Alembert operator is
\[
    \Box=\partial_\mu\partial^\mu.
\]
With our signature and \(c=1\), this is
\[
    \Box=\partial_t^2-\nabla^2.
\]

A relativistic scalar field equation can therefore be written compactly as
\[
    (\Box+m^2)\phi=0.
\]
This is the Klein--Gordon equation. Its compact form displays its Lorentz
covariance: \(\Box\) is a Lorentz-invariant differential operator, and \(m^2\) is
a scalar parameter.

\section{Covariant and Non-Covariant Forms of Equations}

A covariant equation is written in a way that makes its transformation
properties clear. A non-covariant equation may describe the same physics, but it
hides the spacetime structure by separating time and space.

For example, the Klein--Gordon equation can be written covariantly as
\[
    (\Box+m^2)\phi=0.
\]
In components, using \(\Box=\partial_t^2-\nabla^2\), the same equation is
\[
    \partial_t^2\phi-\nabla^2\phi+m^2\phi=0.
\]
The second form is useful for solving initial-value problems, but the first form
makes Lorentz covariance manifest.

Maxwell's equations provide an even stronger example. In vector notation they
are
\[
\begin{aligned}
    \nabla\cdot\vect{E} &= \rho,\\
    \nabla\cdot\vect{B} &= 0,\\
    \nabla\times\vect{E} + \partial_t\vect{B} &= 0,\\
    \nabla\times\vect{B} - \partial_t\vect{E} &= \vect{J},
\end{aligned}
\]
in units adapted to the chosen conventions. In covariant notation these become
\[
    \partial_\mu F^{\mu\nu}=J^\nu,
    \qquad
    \partial_{[\lambda}F_{\mu\nu]}=0.
\]
The second equation can also be written as
\[
    dF=0
\]
in differential-form notation.

\begin{examplebox}
A non-covariant form is often better for calculation in a chosen frame. A
covariant form is better for understanding the structure of the theory. In good
relativistic field theory, both descriptions are available and equivalent.
\end{examplebox}

Covariance does not mean that all components look numerically the same in every
frame. It means that the equation as a geometric statement is frame-independent.
The components change, but the tensorial relation remains valid.

\section{Invariant Actions}

The action principle is the central organizing principle of this course. In a
relativistic field theory, the action should be a Lorentz scalar. This means
that its value should not depend on the inertial coordinate system used to
compute it.

For a scalar field, a standard action is
\[
    S[\phi]=\int \mathcal{L}\,d^4x,
\]
with
\[
    \mathcal{L}
    = \frac{1}{2}\partial_\mu\phi\partial^\mu\phi
      -\frac{1}{2}m^2\phi^2.
\]
The term \(\partial_\mu\phi\partial^\mu\phi\) is a Lorentz scalar because the
index \(\mu\) is contracted. The term \(\phi^2\) is also a scalar. Therefore the
Lagrangian density is a scalar, and the action is invariant under Lorentz
transformations.

More generally, relativistic Lagrangian densities are built by contracting all
spacetime indices. For electromagnetism, the basic invariant is
\[
    F_{\mu\nu}F^{\mu\nu}.
\]
The Maxwell action in vacuum is proportional to
\[
    S[A]= -\frac{1}{4}\int F_{\mu\nu}F^{\mu\nu}\,d^4x.
\]
The factor \(-1/4\) is conventional, but useful: it gives standard normalizations
for the equations of motion and energy density.

\begin{definitionbox}
A relativistic action is \emph{Lorentz invariant} if it is unchanged under
Lorentz transformations. In practice, this is achieved by building the
Lagrangian density from Lorentz scalars and integrating over the invariant
spacetime volume element.
\end{definitionbox}

In flat spacetime, \(d^4x\) is invariant under proper Lorentz transformations
because \(|\det\Lambda|=1\). In curved spacetime, the invariant volume element
is no longer simply \(d^4x\), but \(\sqrt{|g|}\,d^4x\), where \(g\) is the
determinant of the metric. This generalization will become important when field
theory is coupled to gravity.

\section{Why Relativistic Field Theory Is Natural}

Relativistic field theory is not just classical mechanics rewritten with more
indices. It is the natural framework for local physics compatible with special
relativity.

First, special relativity treats space and time as parts of one structure.
Therefore a local dynamical variable should naturally be a function on
spacetime:
\[
    \phi=\phi(x).
\]
Second, causality restricts how disturbances may propagate. Relativistic field
equations must respect the light-cone structure of Minkowski spacetime. Third,
energy and momentum form a four-vector, and conservation laws are naturally
written as divergence equations such as
\[
    \partial_\mu T^{\mu\nu}=0.
\]

Fourth, interactions that are local in spacetime are naturally described by
local Lagrangian densities. A typical action has the form
\[
    S=\int \mathcal{L}(\phi^A,\partial_\mu\phi^A,x)\,d^4x.
\]
This expression says that the total action is assembled from local spacetime
contributions. That locality is one of the deepest structural assumptions of
classical and quantum field theory.

Finally, the same language prepares the transition to quantum field theory. The
objects that will later be quantized are not particle trajectories, but fields,
actions, symmetries, currents, and propagators. Classical relativistic field
theory is therefore the structural bridge between classical mechanics and
quantum field theory.

\begin{summarybox}
Special relativity turns field theory into spacetime geometry. Fields are
functions or tensorial objects on Minkowski spacetime; equations should be
Lorentz covariant; actions should be Lorentz invariant; and causal propagation
is organized by the light cone. This is the natural language for scalar fields,
electromagnetism, Yang--Mills theory, spinor fields, and eventually gravity.
\end{summarybox}
